
Monte Carlo simulations are regarded as the most precise technique for neutronic calculations due to their capacity to employ continuous angle and energy variables and accurately model geometries. Nevertheless, these simulations incur substantial computational expense because they rely on stochastic simulations of individual particles, particularly when striving for acceptable uncertainty levels in localized measurements like power distribution. This computational burden renders Monte Carlo methods less practical for applications demanding a sequence of neutronic calculations, such as transient simulations or iterative processes for thermal-hydraulic feedback.

To overcome this, alternative approaches have been proposed involving the pre-calculation of fission matrices at various reference states using Monte Carlo simulations. These pre-calculated matrices can then be leveraged to rapidly execute neutronic calculations for different reactor conditions. A fission matrix, defined as an n×n matrix, spatially discretizes the reactor geometry into $N_c$ cells. The entry $a_{ij}$ within this matrix signifies the expected number of neutrons generated in cell '$i$' per neutron originating from cell '$j$'. In multiplying systems, the principal eigenvector $F$ of the fission matrix $\textbf{A}$ represents the fission source distribution, while its principal eigenvalue corresponds to the multiplication factor $k_{eff}$.
Fission matrices have been suggested as a means to eliminate inactive cycles in Monte Carlo calculations and to achieve unbiased solutions in problems characterized by high dominance ratios. Their utility for swift neutronic calculations has been extended to various applications, including spent-fuel pools, whole-core power reactor calculations, transient simulations of novel reactor designs, nuclear thermal propulsion reactors, and research reactors. The current investigation builds upon prior research by applying fission matrix methods to the VTB Heat Pipe Micro Reactor (HPMR). The model is based on the Micro-Reactor design under analysis by the NEAMS Micro-Reactor Application Drivers area (REFERENCE STAUFF). As a modeling exercise, a micro-reactor concept was designed at ANL to gather some of the most pressing modeling challenges faced by the micro-reactor industry, which are:  The use of heat pipe technologies to remove the nuclear heat, the use of TRi-structural ISOtropic (TRISO) fuel to enable operations at very high temperatures and the use of rotating control rod drums in the radial reflectors. A more detailed description of the model is given in \Cref{section:VTB HPMR model} The study additionally quantifies the computational time savings realized by employing the fission matrix method compared to full Monte Carlo simulations in this iterative procedure.
